\documentclass[10pt,a4paper]{amsart}
\usepackage[margin=19mm]{geometry}
\usepackage{amsmath,amssymb,mathtools}
\usepackage[scr=rsfs,cal=euler]{mathalfa}
\usepackage{xfrac}
\usepackage[unicode,hidelinks,bookmarks=false]{hyperref}
\newcommand{\ie}{i.e.\ }
\newcommand{\cf}{cf.\ }
\newcommand{\resp}{respectively\ }
\newcommand{\Subsection}{\subsection}
\providecommand{\nobreakdash}{\nobreak\hbox{-}}
\setlength{\emergencystretch}{2em}
\multlinegap0pt
\usepackage{bm}
\usepackage{bbm}
\usepackage{dsfont}
\usepackage{xspace}
\usepackage{cancel}
\usepackage{mathtools}
\usepackage{amsthm}
\makeatletter
\@ifundefined{thm}{\newtheorem{thm}{Theorem}}{}
\makeatother
\title[J-contractive operator valued functions, vector valued de Br]{J-contractive operator valued functions, vector valued de Branges spaces and functional models}
\author{Bharti Garg, Santanu Sarkar}
\date{Source: arXiv:2411.01921v2 (CC BY 4.0)}
\begin{document}
\maketitle

\noindent\emph{Source and licence.} J-contractive operator valued functions, vector valued de Branges spaces and functional models. Version arXiv:2411.01921v2 (CC BY 4.0), distributed under the Creative Commons Attribution 4.0 International licence, as recorded in the arXiv licence metadata for this exact version and shown on the abstract page https://arxiv.org/abs/2411.01921v2. Full rights notice: licenses/arxiv-2411.01921-CC-BY-4.0.txt.\par\smallskip

\noindent\textit{Editorial scope.} Counted unit: the Thm whose source label is \texttt{Theorem 3.3} in the article (source0.tex lines 318--320, source label \texttt{Theorem 3.3}), delivered together with the article's own complete original proof of it (source0.tex lines 321--463, one \texttt{proof} environment of 143 lines). No mathematics is written, repaired or re-derived: statement and proof are the article's own text line for line. Required context retained uncounted: Equation 2.2 (lines 152--154); Theorem 3.1 (lines 193--202). Every definition, display and label of those lines is retained unchanged. Omitted from this package and not redistributed: 1--151, 155--192, 203--317, 464--993. Nothing inside a retained excerpt is omitted or altered: every retained body is delivered exactly as the article's lines are written, so no omission receipt of any kind accompanies this package. Citations: the retained excerpts carry no citation command at all, so no bibliography entry of the article is transcribed and no reference is redistributed. Reference adaptation: none was needed and none was made; every reference inside the delivered unit resolves inside it, so no unresolved reference or citation is produced. Printed slips kept literal: no printed word, symbol, hypothesis or number of the counted statement or of its proof was altered. Layout and class: the article's own class is \texttt{amsart} with packages including times, mathtools, appendix, bm, pifont, amssymb, latexsym, amsmath, amsfonts, eucal, xcolor; the wrapper is the lane's pinned \texttt{amsart} editorial wrapper at 10pt on A4, which restates the theorem-like environments the retained text needs and adds the article's own macro definitions that the retained text needs (none), with extra wrapper packages (bm, bbm, dsfont, xspace, cancel, mathtools). The delivered numbering is the unit's own. Assets: no retained span contains a figure environment, a graphics-inclusion command, a PDF-page inclusion, a table or a TikZ picture; no image file is copied into this package, the package declares no asset dependency and no image file is redistributed with it. Licence: version arXiv:2411.01921v2 of this article is distributed under the Creative Commons Attribution 4.0 International licence, as recorded in the arXiv licence metadata for this exact version and shown on the abstract page https://arxiv.org/abs/2411.01921v2. A full rights notice is delivered at licenses/arxiv-2411.01921-CC-BY-4.0.txt. Characters: every retained excerpt is pure ASCII, so the delivered engine, XeLaTeX with its default Latin Modern text font, reports no missing character.
\begin{equation}\label{Equation 2.2}
U(\lambda)JU(\bar{\lambda})^* = J = U(\bar{\lambda})^*JU(\lambda)~ \mbox{for all}~ \lambda \in \Omega.
\end{equation}

\begin{thm}\label{Theorem 3.1}
Let $\mathcal{H}$ be an RKHS of vector valued analytic functions defined on $\Omega$ which is symmetric about the real axis. Then $\mathcal{H}$ is a $\mathcal{H}(U)$ space if and only if the following conditions hold:
\begin{enumerate}
\item $\mathcal{H}$ is $R_{\alpha}$ invariant for all $\alpha \in \Omega$.
\item The de Branges identity given by  \begin{equation}
\langle R_{\alpha}f,g \rangle -\langle f,R_{\beta}g \rangle -(\alpha-\bar{\beta})\langle R_{\alpha}f, R_{\beta}g \rangle= 2\pi i g(\beta)^* Jf(\alpha),
\end{equation}
holds true for all $f,~g \in \mathcal{H}$ and $\alpha,~\beta \in \Omega$.
\end{enumerate}  
\end{thm}

\begin{thm} \label{Theorem 3.3}
Consider a simple, closed, densely defined, symmetric operator $E$ on the Hilbert space $\mathfrak{X}$ with infinite deficiency indices. Let $\mathcal{H}_1$ denote the RKHS constructed as described above. Then $\mathcal{H}_1$ is a $\mathcal{H}(U)$ space, where $U(\lambda)$ is an analytic operator valued function satisfying $\eqref{Equation 2.2}$ on $\Omega$, with the signature operator $J = J_1$.
\end{thm}

\begin{proof}
In order to demonstrate that $\mathcal{H}_1$ is a $\mathcal{H}(U)$ space, we employ the characterization as provided by Theorem $\ref{Theorem 3.1}$. First, we establish that $\mathcal{H}_1$ is $R_\alpha$ invariant for every $\alpha \in \Omega$. Consider the action of $R_\alpha$ on a function in $\mathcal{H}_1$:
\begin{eqnarray}
\left(R_{\alpha} \begin{bmatrix} 
    f_Y  \\
    \tilde{f}_Y \\
    \end{bmatrix}\right) (\xi)\nonumber
&=&  \begin{bmatrix} 
    (R_{\alpha}f_Y)(\xi)  \\
    (R_{\alpha}\tilde{f}_Y)(\xi) \\ 
    \end{bmatrix}\\ \nonumber
\vspace{0.2cm}       
&=&  \begin{bmatrix} 
    \frac{f_Y(\xi)-f_Y(\alpha)}{\xi-\alpha}  \\
    \frac{\tilde{f}_Y(\xi)- \tilde{f}_Y(\alpha)}  {\xi-\alpha}\\
    \end{bmatrix}\\    
\vspace{0.2cm}       
&=&  \begin{bmatrix} \label{Equation 3.20}
    \frac{P_Y(\xi)f-P_Y(\alpha)f}{\xi-\alpha}  \\
    \frac{\tau_Y(\xi)f- \tau_Y(\alpha)f}{\pi      i(\xi-\alpha)}\\
    \end{bmatrix}.
\end{eqnarray}
Now, let $g$ and $h$ be such that 
$$
f
=(E-\xi I)g+P_Y(\xi)f
=(E-\alpha I)h+P_Y(\alpha)f.
$$
This implies $$ h=\frac{P_Y(\xi)f - P_Y(\alpha)f}{\xi-\alpha}+(E-\xi I)\left( \frac{g-h}{\xi-\alpha}\right).$$
From $\eqref{Equation 3.20}$, we obtain $$  \left(R_{\alpha} \begin{bmatrix} 
    f_Y  \\
    \tilde{f}_Y \\
    \end{bmatrix}\right) (\xi)= \begin{bmatrix} 
    P_Y(\xi)h  \\
    \frac{\tau_Y(\xi)h}{\pi i} \\
    \end{bmatrix}    =  \begin{bmatrix} 
    h_Y  \\
    \tilde{h}_Y \\
    \end{bmatrix} (\xi).$$
Next, we aim to demonstrate the validity of the de Branges identity for all $\alpha, \beta \in \Omega$ and $\begin{bmatrix} 
f_Y  \\
\tilde{f}_Y \\
\end{bmatrix},  \begin{bmatrix} 
g_Y  \\
\tilde{g}_Y \\
\end{bmatrix} \in \mathcal{H}_1$: \begin{multline} \label{Equation 3.21}
\left\langle R_\alpha  \begin{bmatrix} 
	f_Y  \\
	\tilde{f}_Y\\
	\end{bmatrix},  \begin{bmatrix} 
	g_Y  \\
	\tilde{g}_Y \\
	\end{bmatrix} \right\rangle  -\left\langle  \begin{bmatrix} 
	f_Y  \\
	\tilde{f}_Y \\
	\end{bmatrix}, R_{\beta} \begin{bmatrix} 
	g_Y  \\
	\tilde{g}_Y \\
	\end{bmatrix} \right\rangle - (\alpha- \overline{\beta}) \left\langle R_{\alpha}\begin{bmatrix} 
	f_Y  \\
	\tilde{f}_Y \\
	\end{bmatrix},  R_{\beta}\begin{bmatrix} 
	g_Y  \\
	\tilde{g}_Y \\
	\end{bmatrix} \right\rangle 
	\\
	= 2\pi i \begin{bmatrix} 
	g_Y  \\
	\tilde{g}_Y \\
	\end{bmatrix}^*(\beta) J \begin{bmatrix} 
	f_Y  \\
	\tilde{f}_Y\\
	\end{bmatrix}(\alpha).
	\end{multline}
The left-hand side (L.H.S.) of $\eqref{Equation 3.21}$ is expressed as
\begin{multline} \label{Equation 3.22}
\left\langle   \begin{bmatrix} 
	h_Y  \\
	\tilde{h}_Y \\
	\end{bmatrix},  \begin{bmatrix} 
	g_Y  \\
	\tilde{g}_Y \\
	\end{bmatrix} \right\rangle  -\left\langle  \begin{bmatrix} 
	f_Y  \\
	\tilde{f}_Y \\
	\end{bmatrix}, \begin{bmatrix} 
	l_Y  \\
	\tilde{l}_Y \\
	\end{bmatrix} \right\rangle - (\alpha- \overline{\beta}) \left\langle \begin{bmatrix} 
	h_Y  \\
	\tilde{h}_Y \\
	\end{bmatrix},  \begin{bmatrix} 
	l_Y  \\
	\tilde{l}_Y \\
	\end{bmatrix} \right\rangle 
\\
= 2\langle h,g \rangle	-2\langle f,l \rangle-2 (\alpha- \overline{\beta})\langle h, l \rangle	
\end{multline}	
where $h$ and $l$ are defined as 
\begin{equation} \label{Equation 3.23}
f=(E-\alpha I)h +P_Y(\alpha)f \hspace{0.3cm}\text{and} \hspace{0.3cm}
g=(E-\beta I)l+P_Y(\beta)g.
\end{equation}
Observe the following two equalities: $$ 2\langle h,g\rangle =2 \langle h,g-P_Y(\beta)g \rangle+ 2\langle h, P_Y(\beta)g \rangle$$ and \begin{eqnarray*}
\langle h,g-P_Y(\beta)g \rangle &=& \langle h,(E- \beta I)l \rangle \\
&=& \langle Eh,l \rangle- \langle \overline{\beta} h,l\rangle\\
&=& \langle (E-\alpha I)h,l \rangle- \langle (\overline{\beta}-\alpha) h,l\rangle 
\end{eqnarray*} Substituting these into $\eqref{Equation 3.22}$, we obtain $$\text{L.H.S. of $(\ref{Equation 3.21})$}=2\langle h,P_Y(\beta)g \rangle - 2 \langle P_Y(\alpha)f,l \rangle.$$ 
Evaluating R.H.S. of $(\ref{Equation 3.21})$ with $J= \begin{bmatrix}
0 & I \\ 
I & 0
\end{bmatrix}$, we get 
\begin{eqnarray*}
2\pi i  \begin{bmatrix} 
	g_Y  \\
	\tilde{g}_Y\\
	\end{bmatrix}^*(\beta) J  \begin{bmatrix} 
	f_Y  \\
	\tilde{f}_Y \\
	\end{bmatrix}(\alpha) 
&=& 2\pi i  \left\langle J \begin{bmatrix} 
	f_Y  \\
	\tilde{f}_Y \\
	\end{bmatrix}(\alpha) ,  \begin{bmatrix} 
	g_Y  \\
	\tilde{g}_Y \\
	\end{bmatrix}(\beta) \right\rangle \\
&=&	 2 \pi i  \left\langle  \begin{bmatrix}
0 & I \\ 
I & 0
\end{bmatrix} 
\begin{bmatrix} 
	P_Y(\alpha)f  \\
	\frac{\tau_Y(\alpha)f}{\pi i} \\
	\end{bmatrix} ,  \begin{bmatrix} 
	P_Y(\beta)g  \\
	\frac{\tau_Y(\beta)g}{\pi i} \\
	\end{bmatrix} \right\rangle \\
&=& 2\langle \tau_Y(\alpha)f,P_Y(\beta)g \rangle - 2 \langle P_Y(\alpha)f,\tau_Y(\beta)g \rangle\\
&=& 2\langle h,P_Y(\beta)g \rangle - 2 \langle P_Y(\alpha)f,l \rangle.
\end{eqnarray*}
Hence, the result is proved.
\end{proof}

\end{document}
